% 原文 pp.1–2（印刷 pp.1119–1120）
\subsection*{导言}

我这学期讲座的目标，是给出规范场论动力学的一个导引。我们大体上不再处理诸如“什么是量子场论”“哪些量子场论存在”这类形式问题；相反，给定一个量子场论，我们想知道它如何行为——一个给定理论最有意思的物理性质与数学性质是什么；以及通过研究它，我们能学到什么。

从形式问题转向动力学，在数学与物理的许多领域都有类似的情形。例如，面对力学中的 Newton 方程，或流体的 Navier--Stokes 方程，问题之一是要证明该问题是适定的：对给定的初始条件，存在唯一的解。一旦确立了这一点，人们接着就想知道一个给定系统如何行为——这本质上意味着，它在长时间下如何行为。毕竟，有一个适定的微分方程支配着系统，只意味着人们知道短时间的行为；当人们想要解出方程、判定大时间尺度上会发生什么时，才会浮现出一类性质完全不同的深刻问题。

在量子场论中，对应的情形是：要确立一个给定理论存在，主要是一件控制短距离行为的事情，所用工具例如重整化理论，而在微妙的场合则是渐近自由。由于我们在相对论框架下工作，“短距离”既意味着短的空间距离，也意味着短时间。在控制了短距离行为、并使自己相信一个给定理论存在之后，我们接着想要尽可能完整地“求解”这个理论，也就是尽我们所能判定大距离上会发生什么。

这样的研究，要求一种与我们迄今所发展的思维方式颇为不同的思维方式。存在各种类型的大量有趣例子，也存在大量对研究它们最有效的方法。人们必须培养出这样的经验：对给定的模型，提出什么样的问题、采用什么样的方法才最有意义；必须学会分辨哪类信息是重要的、哪类信息是可获得的。看来，只有通过在一些例子上反复练习，才能获得这种经验。在为这学期的讲座挑选例子时，我尽量选择那些相对简单、从物理观点看富有教益、并且具有一定数学意义的例子。此外，我把一个目标当作边界条件：最后要以对四维 $\mathcal{N}=2$ 超对称 Yang--Mills 理论行为的解释收尾（这与四维流形的 Donaldson 理论相关）；这一目标有助于组织部分选题。

关于超对称的习题（见 \texttt{[I-Homework]}）是作为某些讲座的背景材料给出的；其余习题集则汇集于本书之末。
